% Generated draft: requires mathematical and source review.
\documentclass[11pt]{article}
\usepackage[margin=25mm]{geometry}
\usepackage{fontspec}
\setmainfont{Latin Modern Roman}
\usepackage{amsmath,amssymb,mathtools,unicode-math}
\setmathfont{Latin Modern Math}
\usepackage{enumitem,needspace,xurl,hyperref}
\setlength{\parindent}{0pt}
\setlength{\parskip}{5pt}
\setlength{\emergencystretch}{3em}
\newcommand{\N}{\mathbb{N}}
\newcommand{\Z}{\mathbb{Z}}
\newcommand{\Q}{\mathbb{Q}}
\newcommand{\R}{\mathbb{R}}
\newcommand{\C}{\mathbb{C}}
\newcommand{\F}{\mathbb{F}}
\newcommand{\A}{\mathbb{A}}
\newcommand{\PP}{\mathbb{P}}
\newcommand{\E}{\mathbb{E}}
\newcommand{\eps}{\varepsilon}
\newcommand{\dd}{\,\mathrm{d}}
\begin{document}
\section*{UnsolvedMath ID: 700620}
\textbf{Source problem code:} catalogue entry
\subsection{Definitions and mathematical statement}
For a partially ordered collection \((\mathfrak M,\preceq)\), an element \(\mathcal V\in\mathfrak M\) is minimal if
\[
\forall\mathcal W\in\mathfrak M,\qquad
\mathcal W\preceq\mathcal V\Longrightarrow\mathcal W=\mathcal V.
\]
This defines abstract minimality only; the supplied record does not specify the collection and order intended for mixed varieties, or whether any trivial objects are to be excluded.

\textbf{Problem (V. S. Anashin, Kourovka Notebook 10.3).} Characterize minimal mixed varieties of groups, either in terms of bases of mixed identities or in terms of generating groups.

\textbf{Unresolved formulation.} The supplied record does not define a mixed identity, a mixed variety, a basis of mixed identities, or generation of a mixed variety by groups. In particular, it does not specify whether identities involve coefficients or other additional data, or what structures and maps would preserve such data. These notions must be supplied from the intended source before the requested characterization can be stated as a fully specified necessary-and-sufficient condition. The abstract definition above does not resolve these omissions or assert a particular interpretation.
\subsection{Short English statement}
Anashin asks for a characterization of minimal mixed varieties of groups using either bases of mixed identities or generating groups. The supplied record omits the definitions needed to specify this target precisely. Its status is recorded as open; that label has not been independently verified.
\subsection{Sources}
\begin{itemize}
\item Local UnsolvedMath registry, permanent ID \(700620\), exact source problem code \texttt{catalogue entry}, title \emph{Kourovka Notebook Problem 10.3}. The supplied record labels the problem open.
\item The supplied attribution identifies V. S. Anashin as proposer, with proposed year 1986, and cites \emph{The Kourovka Notebook (Unsolved Problems in Group Theory)}, 10th issue (1986), problem 10.3, as reproduced in the 21st issue edition (2026), \texttt{arXiv:1401.0300v48}. The collection description names E. I. Khukhro and V. D. Mazurov as editors. This source attribution is supplied registry information and has not been independently checked.
\end{itemize}
\end{document}
